\section*{第 \ref{ch:b2:powerseries} 章——幂级数}

\begin{solution}{exo:b2:powerseries:1}
$\sum \frac{n^2}{2^n}z^n$：比值 $\frac{(n+1)^2}{2^{n+1}}\cdot
\frac{2^n}{n^2} \to \frac12$：$R = 2$。

$\sum \frac{z^n}{\binom{2n}{n}}$：$\binom{2n}{n} \sim
\frac{4^n}{\sqrt{\pi n}}$（\cref{ex:b2:comparison:centralbinomial}），
故 $\abs{a_n}^{-1} \approx 4^n$（相差多项式因子）：$R = 4$
（比值判别法：$\frac{\binom{2n}{n}}{\binom{2n+2}{n+1}} =
\frac{(n+1)^2}{(2n+1)(2n+2)} \to \frac14$）。

$\sum z^{n!}$：系数 $a_k = 1$（若 $k = n!$），否则为 $0$。对于
$\abs z < 1$ 时，$\sum \abs z^{n!}$ 收敛（由几何级数控制）；
当 $\abs z \geq 1$ 时，通项不趋于 $0$：$R = 1$。

$\sum (2 + (-1)^n)^n z^n$：系数为 $3^n$（$n$ 为偶数）与 $1$
（$n$ 为奇数）。要使 $a_nr^n$ 有界，需 $3r \leq 1$；而 $r <
\frac13$ 即可：$R = \frac13$。
\end{solution}

\begin{solution}{exo:b2:powerseries:2}
三者的\omterm{def:b2:powerseries:radius}{收敛半径}均为 $1$（比值判别法）。在 $z = 1$ 处：$\sum 1$
发散；$\sum\frac1n$ 发散；$\sum\frac{1}{n^2}$ 收敛。在
$z = -1$ 处：$\sum(-1)^n$ 发散；$\sum\frac{(-1)^n}{n}$ 收敛
（交错级数）；$\sum\frac{(-1)^n}{n^2}$ \omterm{def:b2:series:def}{绝对收敛}。
\omterm{def:b2:powerseries:radius}{收敛半径}看不出边界处的情形。
\end{solution}

\begin{solution}{exo:b2:powerseries:3}
由 $\frac{1}{1-x} = \sum x^n$，逐项求导并乘以 $x$
（\cref{thm:b2:powerseries:calculus}）：
\[
\sum n x^n = \frac{x}{(1-x)^2} .
\]
再求一次导，再乘以 $x$：
\[
\sum n^2 x^n = x\,\frac{\dd}{\dd x}\Bigl(\frac{x}{(1-x)^2}\Bigr)
= \frac{x(1 + x)}{(1-x)^3} .
\]
第三个和式是 $-\ln(1 - x)$ 的奇次项部分：
\[
\sum_{n\geq0} \frac{x^{2n+1}}{2n+1}
= \frac{-\ln(1-x) + \ln(1+x)}{2}
= \frac12 \ln\frac{1+x}{1-x}
= \operatorname{artanh} x .
\]
\end{solution}

\begin{solution}{exo:b2:powerseries:4}
部分分式：$\frac{1}{(1-x)(2-x)} = \frac{1}{1-x} -
\frac{1}{2 - x} = \sum x^n - \frac12\sum \bigl(\frac x2\bigr)^n$：
\[
\frac{1}{(1-x)(2-x)} = \sum_{n\geq0} \Bigl(1 -
\frac{1}{2^{n+1}}\Bigr)x^n,
\qquad R = 1 .
\]

$\ln(1 + x + x^2) = \ln\frac{1 - x^3}{1 - x} = \ln(1 - x^3) -
\ln(1 - x) = \sum_{n\geq1}\frac{x^n}{n} -
\sum_{m\geq1}\frac{x^{3m}}{m}$：$x^n$ 的系数为
$\frac1n$（$3 \nmid n$ 时），而 $\frac1n - \frac{3}{n} = -\frac2n$（
$3 \mid n$ 时）。\omterm{def:b2:powerseries:radius}{收敛半径}为 $1$（最近的奇点来自
$\ln(1-x^3)$ 的级数）。
\end{solution}

\begin{solution}{exo:b2:powerseries:5}
$\sum_{m \geq 0} x^m$（系数为 $1$）与
$\sum_{k\geq1} \frac{x^k}{k}$（系数为 $\frac1k$，$k \geq
1$）的\omterm{thm:b2:series:fubini}{柯西乘积}，二者在 $\abs x < 1$ 时\omterm{def:b2:series:def}{绝对收敛}，其乘积中
$x^n$ 的系数为 $\sum_{k=1}^{n} \frac1k \cdot 1 = H_n$。
由此
\[
\Bigl(\sum x^m\Bigr)\Bigl(\sum \frac{x^k}{k}\Bigr)
= \frac{1}{1-x}\cdot\bigl(-\ln(1-x)\bigr)
= \sum_{n\geq1} H_n x^n .
\]
\end{solution}

\begin{solution}{exo:b2:powerseries:6}
将递推式两边乘以 $x^{n+1}$ 并求和（$\abs x <
\frac12$）：
\[
U(x) - 1 = 2x\,U(x) + \sum_{n\geq0} n\,x^{n+1}
= 2x\,U(x) + \frac{x^2}{(1-x)^2} ,
\]
利用\cref{exo:b2:powerseries:3}，于是
\[
U(x) = \frac{1}{1 - 2x}\Bigl(1 + \frac{x^2}{(1-x)^2}\Bigr)
= \frac{1 - 2x + 2x^2}{(1-2x)(1-x)^2} .
\]
部分分式（用遮盖法，在 $x = \frac12$ 处得系数 $2$；
在二重极点 $x = 1$ 处系数为 $-1$；中间
项系数在 $x = 0$ 处取值即为零）：
\[
U(x) = \frac{2}{1-2x} - \frac{1}{(1 - x)^2} .
\]
将两项分别展开：
\[
u_n = 2\cdot 2^n - (n + 1) = 2^{n+1} - n - 1 .
\]
（验算：$u_0 = 1$，$u_1 = 2u_0 + 0 = 2 = 4 - 2$。）
\end{solution}

\begin{solution}{exo:b2:powerseries:7}
归纳法：$f'(x) = \frac{2}{x^3}\eu^{-1/x^2}$，若 $f^{(n)}(x)
= P_n(\tfrac1x)\eu^{-1/x^2}$，则
\[
f^{(n+1)}(x) =
\Bigl(-\frac{1}{x^2}\,P_n'\Bigl(\frac1x\Bigr) +
\frac{2}{x^3}\,P_n\Bigl(\frac1x\Bigr)\Bigr)\eu^{-1/x^2} :
\]
仍为上述形式。在 $0$ 处，差商 $\frac{f^{(n)}(h)}{h} =
\frac1h P_n(\frac1h)\eu^{-1/h^2} \to 0$，当 $h \to 0$ 时，由于
$Q(u)\,\eu^{-u^2} \to 0$（$u \to \pm\infty$），对任意多项式
$Q$ 均成立（指数衰减胜过任何多项式）：由归纳法，所有 $f^{(n)}(0)$ 存在
且都为零，每个 $f^{(n)}$ 在 $0$ 处由同一极限
可知\omterm{def:b2:metric:continuity}{连续}。故 $f \in C^\infty$，且在 $0$ 处泰勒级数为零；
其和为 $0 \neq f$，故在 $0$ 处不\omterm{def:b2:powerseries:analytic}{解析}。
\end{solution}

\begin{solution}{exo:b2:powerseries:8}
二项级数：$\sqrt{1-4x} = \sum_{k\geq0}
\binom{1/2}{k}(-4x)^k$。当 $k = n + 1 \geq 1$ 时：
\begin{align*}
\binom{1/2}{n+1}(-4)^{n+1}
&= \frac{\frac12\bigl(\frac12 - 1\bigr)\cdots\bigl(\frac12 -
n\bigr)}{(n+1)!}\,(-4)^{n+1}\\
&= \frac{(-1)^n\,1\cdot3\cdots(2n-1)}{2^{n+1}(n+1)!}\,(-4)^{n+1}
= -\frac{2}{n+1}\cdot\frac{(2n)!}{n!\,n!} ,
\end{align*}
利用 $1\cdot3\cdots(2n-1) = \frac{(2n)!}{2^n n!}$。于是
\[
C(x) = \frac{1 - \sqrt{1-4x}}{2x}
= \frac{1}{2x}\sum_{n\geq0}\frac{2}{n+1}\binom{2n}{n}x^{n+1}
= \sum_{n\geq0} \frac{1}{n+1}\binom{2n}{n}\,x^n :
\]
$C_n = \frac{1}{n+1}\binom{2n}{n}$。\omterm{def:b2:powerseries:radius}{收敛半径}为 $\frac14$（二项
级数中的自变量为 $4x$）。渐近性由
\cref{ex:b2:comparison:centralbinomial}可得：
\[
C_n \sim \frac{4^n}{\sqrt{\pi}\; n^{3/2}} .
\]
\end{solution}

\begin{solution}{exo:b2:powerseries:9}
设 $A_n = \sum_{k \leq n} a_k \to A$，由\omterm{thm:b2:series:abel}{阿贝尔求和}得，对
$0 \leq x < 1$，
\[
\sum_{n=0}^{\infty} a_n x^n = (1 - x)\sum_{n=0}^{\infty} A_n x^n
\]
（两边均收敛：$(A_n)$ 有界；恒等式由
$a_n = A_n - A_{n-1}$ 及重新编号可得）。又由 $(1 -
x)\sum x^n = 1$:
\[
\sum_n a_nx^n - A = (1-x)\sum_{n} (A_n - A)x^n .
\]
给定 $\varepsilon$，选 $N$，使得 $\abs{A_n - A} \leq
\varepsilon$ 对 $n > N$ 成立；于是
\[
\Bigl|\sum a_nx^n - A\Bigr|
\leq (1-x)\sum_{n \leq N}\abs{A_n - A} + \varepsilon(1 -
x)\sum_{n > N}x^n
\leq (1-x)\,C_N + \varepsilon ,
\]
令 $x \to 1^-$，得上极限 $\leq \varepsilon$，对所有
$\varepsilon$ 成立。因此径向极限为 $A$。

应用：$\sum \frac{(-1)^{n-1}}{n}$ 收敛（交错级数），
且当 $x < 1$ 时其幂级数之和为 $\ln(1 + x)$。由阿贝尔定理，
其和为 $\ln 2$。类似地，$\sum\frac{(-1)^n}{2n+1}x^{2n+1} =
\arctan x$ 得 $\frac\pi4$（$x = 1$ 时）——本科一年级分册中的积分
证明，如今有了结构性的解释。
\end{solution}

\begin{solution}{exo:b2:powerseries:10}
两级数 $\sum\frac{x^n}{n}$ 与 $\sum\frac{x^n}{n+1}$ 的\omterm{def:b2:powerseries:radius}{收敛半径}均为
$1$，且 $\frac{1}{n(n+1)} = \frac1n - \frac1{n+1}$，故当 $0
< \abs x < 1$ 时：
\[
\sum_{n\geq1}\frac{x^n}{n(n+1)}
= -\ln(1-x) - \frac1x\sum_{n\geq1}\frac{x^{n+1}}{n+1}
= -\ln(1-x) - \frac{-\ln(1-x) - x}{x}
= 1 + \frac{1-x}{x}\ln(1-x) .
\]
\omterm{def:b2:powerseries:radius}{收敛半径}为 $1$；并且 $\norm{x^n/(n(n+1))}_{\infty,\intcc{-1}1} =
\frac{1}{n(n+1)}$ 可和，故在
$\intcc{-1}{1}$ 上\omterm{def:b2:funcseq:series}{正规收敛}，从而和函数在该处\omterm{def:b2:metric:continuity}{连续}。当 $x \to
1^-$ 时，$(1-x)\ln(1-x) \to 0$，闭式趋于 $1$——
与裂项求和的结果一致：$\sum\frac{1}{n(n+1)} =
\lim_N\bigl(1 - \frac{1}{N+1}\bigr) = 1$，即 $x = 1$ 时的值。
\end{solution}

\begin{solution}{exo:b2:powerseries:11}
按不动点集合将 $n!$ 个置换分类：选取
$k$ 个不动点（$\binom nk$ 种方式），并对其余
$n - k$ 个对象作错排，得 $n! = \sum_{k=0}^n\binom nk D_{n-k}$。
两边除以 $n!$：
\[
1 = \sum_{k=0}^{n}\frac{1}{k!}\cdot\frac{D_{n-k}}{(n-k)!} ,
\]
这恰好说明 $\eu^x =
\sum\frac{x^k}{k!}$ 与 $D(x) = \sum D_n\frac{x^n}{n!}$ 的\omterm{thm:b2:series:fubini}{柯西乘积}为
$\sum x^n = \frac{1}{1-x}$。两因子在
$\abs x < 1$ 时\omterm{def:b2:series:def}{绝对收敛}（$D_n \leq n!$，故 $D$ 受几何级数
控制）：乘积恒等式成立
（\cref{prop:b2:powerseries:operations}），且
\[
D(x) = \frac{\eu^{-x}}{1-x} .
\]
考虑 $\eu^{-x} = \sum\frac{(-1)^kx^k}{k!}$ 与
$\sum x^m$ 的\omterm{thm:b2:series:fubini}{柯西乘积}中，$x^n$ 的系数为
$\sum_{k=0}^{n}\frac{(-1)^k}{k!}$，由幂级数
系数的唯一性（\cref{thm:b2:powerseries:calculus}）可得：
\[
\frac{D_n}{n!} = \sum_{k=0}^{n}\frac{(-1)^k}{k!}
\xrightarrow[n\to\infty]{} \eu^{-1} :
\]
约 $37\%$ 的置换都是错排，且这一比例
与 $n$ 无关。
\end{solution}

\begin{solution}{exo:b2:powerseries:12}
二项级数，取 $\alpha = -\frac12$，自变量为 $-4x$：
\[
\binom{-1/2}{n}(-4)^n
= \frac{\bigl(-\frac12\bigr)\bigl(-\frac32\bigr)\cdots
\bigl(-\frac{2n-1}2\bigr)}{n!}(-4)^n
= \frac{1\cdot3\cdots(2n-1)}{2^n\,n!}\,4^n
= \frac{(2n)!}{2^n n!}\cdot\frac{2^n}{n!}
= \binom{2n}{n},
\]
利用 $1\cdot3\cdots(2n-1) = \frac{(2n)!}{2^nn!}$。于是
$(1-4x)^{-1/2} = \sum\binom{2n}nx^n$ 对 $\abs{4x} < 1$ 成立。
将其平方（\omterm{thm:b2:series:fubini}{柯西乘积}，由\omterm{def:b2:series:def}{绝对收敛}性可以保证）
并与 $\frac{1}{1-4x} = \sum 4^nx^n$ 比较：
平方后 $x^n$ 的系数为
$\sum_{k=0}^n\binom{2k}k\binom{2n-2k}{n-k}$，由
系数的唯一性可得：
\[
\sum_{k=0}^{n}\binom{2k}{k}\binom{2n-2k}{n-k} = 4^n .
\]
\end{solution}

\begin{solution}{pb:b2:powerseries:1}
\textbf{1.} 取 $a_n = r_n - r_{n+1}$，由分部求和：
\[
\sum_{n=N}^{M} a_nx^n
= r_Nx^N + \sum_{n=N+1}^{M} r_n\bigl(x^n - x^{n-1}\bigr)
- r_{M+1}x^M .
\]
当 $0 \leq x \leq 1$ 时，增量 $x^{n-1} - x^n$ 非负，
且逐项相消得 $x^N - x^M$；记 $s = \sup_{n\geq
N}\abs{r_n}$：
\[
\Bigl|\sum_{n=N}^{M}a_nx^n\Bigr|
\leq s\bigl(x^N + (x^N - x^M) + x^M\bigr) = 2s\,x^N \leq 2s .
\]

\textbf{2.} 由 $r_n \to 0$，$\sup_{n\geq N}\abs{r_n} \to 0$：
第 1 小题恰为一致柯西准则在
$\intcc{0}{1}$ 上，故 $\sum a_nx^n$ 在其上\omterm{def:b2:funcseq:def}{一致收敛}，且
其和函数\omterm{def:b2:metric:continuity}{连续}（见\cref{thm:b2:funcseq:seriestransfer}）。
$1$ 处的值为 $\sum a_n$，故 $1$ 处的\omterm{def:b2:metric:continuity}{连续}性是
\cref{exo:b2:powerseries:9} 的径向极限。

\textbf{3.} 当 $\abs x < 1$ 时，三个幂级数都
\omterm{def:b2:series:def}{绝对收敛}，且 $\bigl(\sum a_nx^n\bigr)\bigl(\sum b_nx^n\bigr)
= \sum c_nx^n$（见\cref{prop:b2:powerseries:operations}）。
由第 2 小题，各因子及乘积两侧在
$\intcc{0}{1}$ 上\omterm{def:b2:metric:continuity}{连续}（其系数级数
由假设收敛）；令 $x \to 1^-$ 代入恒等式，得 $AB = C$。

\textbf{4.} 此时
\[
\abs{c_n} = \sum_{k=0}^{n}
\frac{1}{\sqrt{(k+1)(n-k+1)}}
\geq \sum_{k=0}^{n}\frac{2}{n+2}
= \frac{2(n+1)}{n+2} \geq 1,
\]
由算术--几何平均不等式：$\sqrt{(k+1)(n-k+1)} \leq \frac{(k+1) +
(n-k+1)}{2} = \frac{n+2}{2}$。级数 $\sum c_n$ 的
通项不趋于 $0$：\omterm{thm:b2:series:fubini}{柯西乘积}发散，尽管
两个因子均收敛（交错级数）。

\textbf{5.} 由\cref{exo:b2:powerseries:5}可得
$\frac{-\ln(1-x)}{1-x} = \sum H_nx^n$（$\abs x < 1$）。
逐项求原函数
（见\cref{thm:b2:powerseries:calculus}~(2)），两边在
$0$ 处均为零：
\[
\frac{\bigl(\ln(1-x)\bigr)^2}{2}
= \sum_{n\geq1}\frac{H_n}{n+1}\,x^{n+1} .
\]
递减性：$(n+2)H_n \geq (n+1)H_{n+1}$ 等价于 $H_n \geq
1$，对 $n \geq 1$ 成立；又 $\frac{H_n}{n+1} \sim \frac{\ln
n}{n} \to 0$：在 $x = -1$ 处，级数由
交错级数判别法收敛。代入 $x \mapsto -x$，并
在 $x = 1$ 处应用第 2 小题：
\[
\frac{(\ln 2)^2}{2}
= \sum_{n\geq1}\frac{H_n}{n+1}(-1)^{n+1},
\]
即所述的值。

\textbf{6.} $f(x) = \sum(-1)^nx^n = \frac{1}{1+x} \to \frac12$
当 $x \to 1^-$ 时，阿贝尔可和于 $\frac12$。但
部分和为 $1, 0, 1, 0, \dots$，发散。

\textbf{7.} 给定 $\varepsilon > 0$，取 $m$ 使 $\abs{u_n}
\leq \varepsilon$（$n > m$）；对 $N \geq m$：
\[
\Bigl|\frac{u_1 + \dots + u_N}{N}\Bigr|
\leq \frac{\abs{u_1} + \dots + \abs{u_m}}{N} +
\varepsilon\,\frac{N - m}{N}
\leq \frac{C_m}{N} + \varepsilon,
\]
于是 $\limsup \leq \varepsilon$ 对每个 $\varepsilon$ 成立，
即均值趋于 $0$。

\textbf{8.} 当 $0 \leq x \leq 1$ 时，$1 - x^n = (1-x)(1 + x +
\dots + x^{n-1}) \leq n(1-x)$，于是
\[
\Bigl|\sum_{n=0}^N a_n(1 - x_N^n)\Bigr|
\leq (1 - x_N)\sum_{n=1}^N n\abs{a_n}
= \frac1N\sum_{n=1}^{N}n\abs{a_n} .
\]
当 $n > N$ 时，$\abs{a_n} = \frac{n\abs{a_n}}{n} \leq
\frac{1}{N}\sup_{m>N}m\abs{a_m}$，且 $\sum_{n>N}x_N^n \leq
\frac{1}{1 - x_N} = N$：
\[
\Bigl|\sum_{n>N}a_nx_N^n\Bigr|
\leq \frac{\sup_{m>N}m\abs{a_m}}{N}\cdot N
= \sup_{m>N}\,m\abs{a_m} .
\]

\textbf{9.} 分解为
\[
A_N - L = \sum_{n=0}^{N}a_n\bigl(1 - x_N^n\bigr)
- \sum_{n>N}a_nx_N^n + \bigl(f(x_N) - L\bigr) .
\]
第一项由第 7 小题知趋于 $0$（
$n\abs{a_n} \to 0$ 的均值），第二项由第 8 小题（上确界趋于
$0$），第三项则因 $x_N \to 1^-$ 且 $f(x) \to L$ 而趋于零。
于是 $A_N \to L$：陶伯定理。

\textbf{10.} 当 $x \in \intco{0}{1}$ 且 $N$ 任意时，
$\sum_{n\leq N}a_nx^n \leq f(x) \leq M$（各项非负）。
令 $x \to 1^-$ 于有限和中，得 $\sum_{n\leq N}a_n \leq M$。
部分和单调递增且有界，故 $\sum a_n$
收敛，再由第 2 小题得 $\lim_{x\to1^-}f(x) =
\sum a_n$。

\textbf{11.} $\sigma_N - L$ 是 $N$ 个数
$A_n - L$（$0 \leq n < N$）的均值，它们趋于 $0$：第 7 小题。

\textbf{12.} $A_n = 1$ 对偶数 $n$，$0$ 对奇数：$A_0 + \dots
+ A_{N-1} = \lceil N/2\rceil$，故 $\sigma_N =
\frac{\lceil N/2\rceil}{N} \to \frac12$，即
第 6 小题的阿贝尔值。

\textbf{13.} 由 $\sigma_N \to L$ 可得 $S_n = O(n)$，从而
$A_n = S_n - S_{n-1} = O(n)$，$a_n = O(n)$：下面各级数的
\omterm{def:b2:powerseries:radius}{收敛半径}都 $\geq 1$。对 $\abs x < 1$，由 $a_n = A_n -
A_{n-1}$ 及 $A_nx^n \to 0$：
\[
(1-x)\sum_n A_nx^n = \sum_n A_nx^n - \sum_n A_nx^{n+1}
= \sum_n a_nx^n = f(x),
\]
同理 $(1-x)\sum S_nx^n = \sum A_nx^n$，故 $f(x) =
(1-x)^2\sum_n S_nx^n = (1-x)^2\sum_n(n+1)\sigma_{n+1}x^n$。
最后 $\sum(n+1)x^n = \frac{1}{(1-x)^2}$
（见\cref{exo:b2:powerseries:3}），即第二个恒等式。

\textbf{14.} 用第一个恒等式减去 $L$ 倍的第二个
恒等式，得：
\[
f(x) - L = (1-x)^2\sum_{n\geq0}(n+1)
\bigl(\sigma_{n+1} - L\bigr)x^n .
\]
给定 $\varepsilon$，取 $N$，使得 $\abs{\sigma_{n+1} - L}
\leq \varepsilon$，当 $n \geq N$ 时；则
\[
\abs{f(x) - L} \leq (1-x)^2 C_N +
\varepsilon(1-x)^2\sum_{n}(n+1)x^n
= (1-x)^2C_N + \varepsilon ,
\]
令 $x \to 1^-$，得 $\limsup \leq \varepsilon$。由此
得 $f(x) \to L$，即弗罗贝尼乌斯定理。

\textbf{15.} $f(x) = \sum(-1)^n(n+1)x^n = \frac{1}{(1+x)^2}$
（对几何级数在 $-x$ 处求导）：阿贝尔值
$\frac14$。部分和：$A_{2k} = k+1$，$A_{2k+1} = -(k+1)$
（由归纳法即得）。于是 $S_{2m-1} = 0$（相邻两项
两两相消），且 $S_{2m} = m + 1$，故
\[
\sigma_{2m} = \frac{S_{2m-1}}{2m} = 0,
\qquad
\sigma_{2m+1} = \frac{m+1}{2m+1} \to \frac12 :
\]
$(\sigma_N)$ 有两个不同的聚点值，因此不是
切萨罗可和的。结合第 6 小题与第 11--14 小题，
层级“收敛 $\Rightarrow$ 切萨罗可和 $\Rightarrow$ 阿贝尔可和”
在两个箭头处均是严格的。

\textbf{16.} 逐项积分所得的级数 $F(x) =
\sum\frac{a_n}{n+1}x^{n+1}$ 与原级数有相同的\omterm{def:b2:powerseries:radius}{收敛半径}，且 $F' = f$ 在
$\intco{0}{1}$ 上成立（\cref{thm:b2:powerseries:calculus}）。因为
$\sum\frac{a_n}{n+1}$ 收敛，由第一部分（第 2 小题）得 $F$
在 $\intcc{0}{1}$ 上\omterm{def:b2:metric:continuity}{连续}。由于 $\int_0^x f = F(x)$（两者
导数相同，且在 $0$ 处的值均为 $0$），
\[
\int_0^x f \xrightarrow[x\to1^-]{} F(1)
= \sum_{n\geq0}\frac{a_n}{n+1} :
\]
因此\omterm{def:b2:integration:improper}{广义积分}存在，且取所述的值。

\textbf{17.} $\frac{\ln(1+x)}{x} =
\sum_{n\geq1}\frac{(-1)^{n-1}}{n}x^{n-1}$（\omterm{def:b2:powerseries:radius}{收敛半径}为 $1$；
在 $0$ 处\omterm{def:b2:metric:continuity}{连续}）。由 $\frac{a_m}{m+1}$ 构成的级数为
$\sum_{n\geq1}\frac{(-1)^{n-1}}{n^2}$，\omterm{def:b2:series:def}{绝对收敛}：
由第 16 小题得 $\int_0^1\frac{\ln(1+x)}{x}\dd x = \eta$。在
\omterm{def:b2:series:def}{绝对收敛}的 $\sum\frac1{n^2}$ 中，按奇偶
分组，得：
\[
\eta = \sum_{\text{奇数}}\frac1{n^2} -
\sum_{\text{偶数}}\frac1{n^2}
= \sum_{n}\frac1{n^2} - 2\sum_{k}\frac1{(2k)^2}
= \Bigl(1 - \frac12\Bigr)\sum_n\frac1{n^2}
= \frac12\sum_{n\geq1}\frac1{n^2} .
\]

\textbf{18.} 由莱布尼茨判别法，$\frac{1}{3n+1}\downarrow0$，级数
收敛。在 $\intco{0}{1}$ 上，$\sum(-1)^nx^{3n} =
\frac{1}{1+x^3}$，且 $\sum\frac{(-1)^n}{3n+1}$ 收敛：
由第 16 小题得 $\sum\frac{(-1)^n}{3n+1} =
\int_0^1\frac{\dd x}{1+x^3}$。部分分式（验证：
$\frac13(x^2-x+1) + \frac{2-x}{3}(1+x) = 1$）：
\[
\int_0^1\frac{\dd x}{1+x^3}
= \frac13\ln2 + \frac13\int_0^1\frac{2-x}{x^2-x+1}\dd x .
\]
将 $2 - x = -\frac12(2x-1) + \frac32$ 改写，则 $\ln(x^2-x+1)$
这一部分在两端为零，而
\[
\frac32\int_0^1\frac{\dd x}{(x-\frac12)^2 + \frac34}
= \frac32\cdot\frac{2}{\sqrt3}
\Bigl[\arctan\frac{2x-1}{\sqrt3}\Bigr]_0^1
= \sqrt3\cdot\frac{\pi}{3} = \frac{\pi}{\sqrt3} .
\]
故总和为 $\frac13\bigl(\ln2 + \frac{\pi}{\sqrt3}\bigr)$。

\textbf{19.} 将 $t = x^2$ 代入
\cref{exo:b2:powerseries:12} 中的级数，并逐项积分
（$(1-x^2)^{-1/2}$ 的原函数中在 $0$ 处取零者是
$\arcsin$）：
\[
\arcsin x = \sum_{n\geq0}
\frac{\binom{2n}n}{4^n(2n+1)}x^{2n+1}
\qquad(\abs x < 1) .
\]
其系数 $\sim \frac{1}{2\sqrt\pi\,n^{3/2}}$
（见\cref{ex:b2:comparison:centralbinomial}），是可和的，故级数
在 $\intcc{-1}{1}$ 上\emph{\omterm{def:b2:funcseq:series}{正规收敛}}，其和
在该区间上\omterm{def:b2:metric:continuity}{连续}，且与\omterm{def:b2:metric:continuity}{连续}的 $\arcsin$
在 $\intoo{-1}{1}$ 上重合，故在 $x = 1$ 处也成立：
\[
\sum_{n\geq0}\frac{\binom{2n}n}{4^n(2n+1)}
= \arcsin 1 = \frac\pi2 .
\]

\textbf{20.} $C_n4^{-n} \sim \frac{1}{\sqrt\pi\,n^{3/2}}$
（\cref{exo:b2:powerseries:8}）：$\sum
C_nx^n$ 在 $\intcc{0}{\frac14}$ 上\omterm{def:b2:funcseq:series}{正规收敛}，故其和在该区间上\omterm{def:b2:metric:continuity}{连续}；
在 $\intoo{0}{\frac14}$ 上，它等于
$\frac{1-\sqrt{1-4x}}{2x}$
（\cref{ex:b2:powerseries:catalan}），其在 $\frac14^-$ 处的极限为
$\frac{1-0}{1/2} = 2$。于是 $\sum_{n\geq0}
C_n4^{-n} = 2$。

\textbf{21.} $f(x) = \sum_{n\geq m}a_nx^n = x^m g(x)$，其中
$g(x) = \sum_{k\geq0}a_{m+k}x^k$；若 $(a_nr^n)$ 有界，则
$(a_{m+k}r^k)$ 亦有界（除以 $r^m$），从而 $g$ 的\omterm{def:b2:powerseries:radius}{收敛半径} $\geq
R$。$g$ \omterm{def:b2:metric:continuity}{连续}，且 $g(0) = a_m \neq 0$，故 $g \neq 0$
在某个 $\intcc{-\delta}{\delta}$ 上成立，且 $f(x) = x^mg(x) \neq 0$
对 $0 < \abs x \leq \delta$ 成立。

\textbf{22.} $d = f - h$ 是 $0$ 附近幂级数之和，
它在非零点 $x_k \to 0$ 处为零。若 $d$ 的某个系数不为零，
则由第 21 题可知 $0$ 的某去心邻域内 $d$ 无零点——
这与 $d(x_k) = 0$ 矛盾。
故 $d$ 的所有系数为零：
$f$ 与 $h$ 系数相同，二者在 $0$ 附近重合。

\textbf{23.} 函数 $h(x) = \frac{1}{1+x^2} =
\sum(-1)^nx^{2n}$（\omterm{def:b2:powerseries:radius}{收敛半径} $1$）满足 $h(\frac1k) =
\frac{1}{1 + 1/k^2} = \frac{k^2}{k^2+1}$。任一取相同值的\omterm{def:b2:powerseries:analytic}{解析函数} $f$
与 $h$ 在 $\frac1k \to 0$ 这些点上一致：
由唯一性定理（第 22 题），$f = \frac{1}{1+x^2}$
在 $0$ 附近成立——即唯一解。

\textbf{24.} 设 $Z$ 为 $I$ 中这样的点的集合：这些点具有一邻域，
在该邻域上 $f$ 恒为零。由定义，此集合为\omterm{def:b2:metric:topology}{开集}，
非空（含子区间）。在 $I$ 中为闭：若 $y
\in I$ 为 $Z$ 中点的极限，则 $y$ 为 $f$ 零点的聚点；
将 $f$ 在 $y$ 处展开为幂级数（由\omterm{def:b2:powerseries:analytic}{解析}性），
并以 $y$ 为中心应用第 21--22 题，
则 $y$ 处的所有系数为零，故 $f \equiv 0$ 在 $y$ 附近：$y
\in Z$。区间是\omterm{def:b2:metric:connected}{连通的}，故 $Z = I$：$f \equiv 0$ 于
$I$ 上。特别地，$\R$ 上在紧集外为零的\omterm{def:b2:powerseries:analytic}{解析函数}
在某区间上为零，从而处处为零：
不存在非零的\omterm{def:b2:powerseries:analytic}{解析}隆起。$C^\infty$ 的情形
则不同：利用
\cref{exo:b2:powerseries:7} 中的平坦函数（例如 $x \mapsto
\eu^{-1/x^2}\mathbf 1_{x>0}$ 及其镜像）粘合，可得光滑的
具紧支集隆起。

\textbf{25.} (i) \omterm{def:b2:funcseq:series}{正规收敛}发生在圆盘内部的紧子圆盘上，
即严格位于收敛圆盘之内；阿贝尔定理将\omterm{def:b2:metric:continuity}{连续}性推广到
边界点，唯一的假设是
系数级数在该点收敛。(ii) 其逆命题
在陶伯条件 $na_n \to 0$（第 9 题）下成立，并且
切萨罗可和性严格介于收敛与
阿贝尔可和之间（弗罗贝尼乌斯，第 14--15 题）。(iii) 给朋友的
一道题：$\sum\frac{(-1)^n}{3n+1} = \int_0^1\frac{\dd x}{1+x^3}$
将几何级数逐项积分至边界，再
作部分分式分解。(iv) 切萨罗平均在傅里叶
级数一章中以费耶尔定理的形式重新出现，其中对部分和取平均
可修复\omterm{def:b2:funcseq:def}{逐点收敛}的失效——同一剂良方，
只是换了病人。
\end{solution}
